\section{Conclusiones}\label{sec:conclusiones}

La ejecución de esta investigación permite concluir que la robustez de un sistema
predictivo en salud depende tanto de la capacidad de los algoritmos como de la
calidad metodológica del flujo de preprocesamiento. Mediante un esquema modular
y reproducible, se observó que los modelos de ensamble, específicamente Random
Forest y Gradient Boosting Machine (GBM), lideraron la capacidad de
discriminación con valores de AUC de hasta 0.949 en validación y 0.942 en prueba
en el escenario clínico-estadístico, y de hasta 0.978 en validación y 0.975 en
prueba en el escenario técnico-estadístico. Estos resultados indican
que, ante variables fisiológicas con alta variabilidad, multicolinealidad y
relaciones no lineales, los algoritmos basados en árboles presentan ventajas
frente a aproximaciones lineales, de margen o basadas en distancia como SVM y
KNN. La regresión logística también mostró un rendimiento relevante, con AUC
superiores a 0.93 en el escenario biológico, lo que sugiere que parte de los
predictores seleccionados mantienen una relación discriminativa estable con la
mortalidad hospitalaria. La evaluación final sobre el conjunto de prueba,
reservada únicamente para el mejor modelo de cada escenario, permitió confirmar
la tendencia observada en validación sin utilizar el test como criterio de
selección.

El aporte principal de este trabajo reside en mostrar que el preprocesamiento
informado por conocimiento clínico no debe evaluarse únicamente por la
optimización métrica, sino también por su capacidad para preservar la validez
metodológica del experimento. La comparación entre la estrategia basada en
lógica clínica y la estrategia técnica reveló una paradoja relevante: aunque los
métodos automatizados pueden alcanzar valores superiores de AUC, dichos
resultados pueden estar influidos por la retención de variables con potenciales
sesgos o riesgo de \textit{Data Leakage}. La eliminación de variables de
seguimiento temporal y de estimaciones derivadas de modelos previos, junto con
la imputación estratificada por grupo diagnóstico, permitió preservar mejor la
coherencia biológica del conjunto de datos. Además, el análisis del carácter
informativo de los valores faltantes en las escalas de dependencia funcional
mostró que la ausencia de datos no siempre representa un error técnico, sino que
puede constituir una señal clínica relevante. Como limitación principal, la
validez de estos hallazgos se encuentra acotada a la cohorte SUPPORT2 y debe ser
interpretada con cautela fuera de este contexto. En conjunto, el estudio
respalda la integración del razonamiento clínico en el preprocesamiento como una
condición necesaria para desarrollar modelos predictivos no solo precisos, sino
también metodológicamente defendibles.
